% Part: lambda-calculus
% Chapter: lambda-definability
% Section: truth-values

\documentclass[../../../include/open-logic-section]{subfiles}

\begin{document}

\olfileid[hi]{lam}{ldf}{tvr}
\olsection{सत्य-मान और संबंध}

शुद्ध लैम्ब्डा कलन में सत्य-मानों को इस प्रकार कूटित कर सकते हैं:
\begin{align*}
  \fn{true} & \ident \lambd[x][\lambd[y][x]]\\
  \fn{false} & \ident \lambd[x][\lambd[y][y]]
\end{align*}

सत्य-मानों को \emph{चयनकों} के रूप में निरूपित किया जाता है, अर्थात् ऐसे
फलनों के रूप में जो दो तर्क स्वीकार करके उनमें से एक लौटाते हैं। सत्य-मान
$\fn{true}$ अपना पहला और $\fn{false}$ अपना दूसरा तर्क चुनता है। उदाहरणार्थ,
$\fn{true}\, M N$ सदा अपचयित होकर $M$ बनता है, जबकि $\fn{false}\, M N$
सदा अपचयित होकर~$N$ बनता है।

\begin{defn}
किसी संबंध $R \subseteq \Nat^n$ को हम !!{lambda definable} कहते हैं यदि
कोई ऐसा पद~$R$ हो कि
\begin{align*}
  R\, \num{n_1} \dots \num{n_k} & \bred \fn{true}
  \intertext{जब भी $R(n_1, \dots, n_k)$ हो, और}
  R\, \num{n_1} \dots \num{n_k} & \bred \fn{false}
\end{align*}
अन्यथा।
\end{defn}

उदाहरणार्थ, संबंध $\fn{IsZero} = \{0\}$, जो $0$ और केवल $0$ के लिए लागू
होता है, निम्न द्वारा !!{lambda definable} है:
\[
  \fn{IsZero} \ident \lambd[n][n (\lambd[x][\fn{false}])\, \fn{true}].
\]
यह कैसे काम करता है? चूँकि चर्च अंक पुनरावर्तकों के रूप में परिभाषित हैं
(ऐसे फलन जो अपने पहले तर्क को दूसरे पर $n$ बार लागू करते हैं), इसलिए हम
आरंभिक मान $\fn{true}$ रखते हैं और पुनरावर्तन के प्रत्येक चरण में पिछले
चरण के परिणाम की परवाह किए बिना $\fn{false}$ लौटाते हैं। यह चरण आरंभिक मान
पर $n$ बार लागू होगा; परिणाम $\fn{true}$ तभी और केवल तभी होगा जब चरण
बिल्कुल लागू न हो, अर्थात् जब $n = 0$।

सत्य-मानों के इस निरूपण के आधार पर हम कुछ सत्य-फलन भी परिभाषित कर सकते
हैं। ये दो, निषेध और संयोजन के निरूपण, हैं:
\begin{align*}
  \fn{Not} & \ident \lambd[x][x\, \fn{false}\, \fn{true}]\\
  \fn{And} & \ident \lambd[x][\lambd[y][xy \,\fn{false}]]
\end{align*}
फलन ``$\fn{Not}$'' एक तर्क स्वीकार करता है और तर्क $\fn{false}$ होने पर
$\fn{true}$ तथा तर्क~$\fn{true}$ होने पर $\fn{false}$ लौटाता है। फलन
``$\fn{And}$'' दो सत्य-मान तर्कों के रूप में स्वीकार करता है और दोनों तर्क
~$\fn{true}$ होने पर ही $\fn{true}$ लौटाना चाहिए। सत्य-मानों को ऊपर वर्णित
चयनकों के रूप में निरूपित किया गया है; अतः जब $x$ कोई सत्य-मान हो और उसे
दो तर्कों पर लागू किया जाए, तब $x$ के $\fn{true}$ होने पर परिणाम पहला तर्क
और अन्यथा दूसरा तर्क होगा। अब $\fn{And}$ अपने दो तर्क $x$ और $y$ लेकर,
अपने पहले तर्क~$x$ को बदले में $y$ और $\fn{false}$ देता है। $x$ को सत्य-मान
मानें तो $x$ के $\fn{true}$ होने पर परिणाम~$y$ और $x$ के $\fn{false}$ होने
पर $\fn{false}$ बनेगा; ठीक यही अपेक्षित है।

ध्यान दें कि यहाँ हम मानते हैं कि $\fn{And}$ को तर्कों के रूप में केवल
सत्य-मान दिए जाते हैं। यदि उसे अन्य पद दिए जाएँ, तो परिणाम (अर्थात्
प्रसामान्य रूप, यदि वह अस्तित्व में हो) सत्य-मान न भी हो सकता है।

\begin{prob}
  सत्य-मानों को $\lambd$-पदों के रूप में कूटित करके समावेशी और अपवर्जी
  वियोजन के सत्य-फलनों को निरूपित करने वाले फलन $\fn{Or}$ और $\fn{Xor}$
  परिभाषित करें।
\end{prob}

\end{document}
